\documentclass{beamer}
\usetheme{Madrid}

\usepackage{amsmath, amssymb, amsthm}
\usepackage{graphicx}
\usepackage{gensymb}
\usepackage[utf8]{inputenc}
\usepackage{hyperref}
\usepackage{tikz}


\title{12.116 Matgeo}
\author{AI25BTECH11012 - Garige Unnathi}
\date{}

\begin{document}

\frame{\titlepage}

% Question frame
\begin{frame}
\frametitle{Question}
Let \textbf{P} and \textbf{Q} be two square matrices of same size. Consider the following statements
\begin{enumerate}
    \item \textbf{P}\textbf{Q} = 0 then \textbf{P} = 0 or \textbf{Q} =0 or both
    \item \textbf{P}\textbf{Q} = \textbf{I} then \textbf{Q}\textbf{P} = \textbf{I}
    \item (\textbf{P} + \textbf{Q})$^2$ = \textbf{P}$^2$ + 2\textbf{P}\textbf{Q} + \textbf{Q}$^2$
    \item (\textbf{P} - \textbf{Q})$^2$ = \textbf{P}$^2$ - 2\textbf{P}\textbf{Q} + \textbf{Q}$^2$
\end{enumerate}
Determine which statements are True and which are False
\end{frame}


% Solution steps
\begin{frame}
\frametitle{Option 1}
This option is False\\
    This can be proven by a simple counter example :\\
    \begin{align}
        \textbf{P} = \begin{bmatrix}1&0\\0&0\end{bmatrix} \quad
        \textbf{Q} = \begin{bmatrix}0&0\\1&0\end{bmatrix}\\
        \textbf{P}\textbf{Q} = \begin{bmatrix}0&0\\0&0\end{bmatrix}
    \end{align}
\end{frame}



\begin{frame}
\frametitle{Option 2}
Given \textbf{P}\textbf{Q} = \textbf{I} \\
     From this we can say that $\lvert P \rvert$ and $\lvert Q \rvert$ is not zero \\
     That is inverse exists
     \begin{align}
         \textbf{Q} = \textbf{P}^{-1}\\
         \textbf{Q}\textbf{P} = \textbf{P}^{-1}\textbf{P} = \textbf{I}
     \end{align}
     Hence this option is true 


\end{frame}

\begin{frame}
\frametitle{Option 3}
\begin{align}
         (\textbf{P}+\textbf{Q})^2 = (\textbf{P}+\textbf{Q})(\textbf{P}+\textbf{Q})\\
          =  \textbf{P}^2 + \textbf{P}\textbf{Q} + \textbf{Q}\textbf{P} + \textbf{Q}^2
     \end{align}
     This equation will not be equal to \textbf{P}$^2$ + 2$\textbf{P}\textbf{Q}$ + \textbf{Q}$^2$\\
     this is because Matrix multiplication is not commutative 
     \begin{align}
         \textbf{P}\textbf{Q} \neq \textbf{Q}\textbf{P}
     \end{align}
     this option is false 

\end{frame}
\begin{frame}
\frametitle{Option 4}
Similarly we can also poove that Option 4 is false\\
Hence only option 2 is true
\end{frame}

\end{document}
